%---------------------------------------------------------

\begin{frame}{Teza}
\alt<1>{
    \textcolor{cfblue}{Twierdzenie} \\
    Niech \(G = (V, E)\) będzie grafem takim, że
    \begin{itemize}
        \item \(G\) jest gładki
    \end{itemize}
    Albo \(G\) zawiera pętlę lub jego rdzeń jest cyklem, albo pozwala wydefiniować dowolną strukturę relacyjną.
}{
    \textcolor{black!40}{
        \textcolor{cfblue}{Twierdzenie} \\
        Niech \(G = (V, E)\) będzie grafem takim, że
        \begin{itemize}
            \item \(G\) jest gładki,
        \end{itemize}
    Albo \(G\) zawiera pętlę lub jego rdzeń jest cyklem, albo pozwala wydefiniować dowolną strukturę relacyjną.
    }
}

\vspace{0.5em}

\only<2>{
    \textcolor{cfblue}{Lemat} \\
    Niech \(G = (V, E)\) będzie grafem takim, że
    \begin{itemize}
        \item \(G\) jest gładki,
        \item \(G\) ma \(\mathrm{al}=1\)
    \end{itemize}
    Albo \(G\) zawiera pętlę, albo pozwala wydefiniować dowolną strukturę relacyjną.
}

\end{frame}

%---------------------------------------------------------

\begin{frame}{Relacja płotkowania}

\begin{columns}[c,onlytextwidth]
  % LABELS
  \begin{column}{0.2\textwidth}
      
  \end{column}
  \begin{column}{0.15\textwidth} 
  \vspace{-3.5em}\\
    $F_1(u,v) \equiv$

    \vspace{4em}

    $F_2(u,v) \equiv$

    \vspace{5em}

    $F_3(u,v) \equiv$
  \end{column}

  \begin{column}{0.82\textwidth}
    \begin{tikzpicture}[
    scale=0.5,
    line width=1pt,
    every node/.style={
        circle,
        fill=black,
        inner sep=2.4pt
    },
    arr/.style={-{Triangle[length=2mm,width=1.5mm]}},
]

    % \node at (-3,-5.5) {\(F_1\)};
    
    \node[circle, fill=black, label=below:{\(u\)}]
            (p0) at (-4,-4) {};

    \node[circle, fill=black, label=below:{\(v\)}]
            (p1) at (-2,-4) {};

    \node[circle, fill=black]
        (p2) at (-3,-2) {};
    

    \draw[arr] (p0) -- (p2);
    \draw[arr] (p1) -- (p2);
\end{tikzpicture}

    \begin{tikzpicture}[
    scale=0.5,
    line width=1pt,
    every node/.style={
        circle,
        fill=black,
        inner sep=2.4pt
    },
    arr/.style={-{Triangle[length=2mm,width=1.5mm]}},
]
   % \node at (-4,-2) {\(F_2\)};
            
    \node[circle, fill=black, label=below:{\(u\)}]
            (q0) at (4,-4) {};
    \node[circle, fill=black, label=below:{}]
            (q1) at (8,-4) {};
    \node[circle, fill=black, label=below:{\(v\)}]
            (q2) at (12,-4) {};
    \node[circle, fill=black, label=below:{}]
            (q3) at (5,-2) {};
    \node[circle, fill=black, label=below:{}]
            (q4) at (7,-2) {};
    \node[circle, fill=black, label=below:{}]
            (q5) at (9,-2) {};
    \node[circle, fill=black, label=below:{}]
            (q6) at (11,-2) {};
    \node[circle, fill=black, label=below:{}]
        (q7) at (6,0) {};
    \node[circle, fill=black, label=below:{}]
        (q8) at (10,0) {};   
        
    \draw[arr] (q0) -- (q3);
    \draw[arr] (q3) -- (q7);
    \draw[arr] (q4) -- (q7);
    \draw[arr] (q1) -- (q4);
    \draw[arr] (q1) -- (q5);
    \draw[arr] (q2) -- (q6);
    \draw[arr] (q5) -- (q8);
    \draw[arr] (q6) -- (q8);
\end{tikzpicture}

    \begin{tikzpicture}[
    scale=0.35,
    line width=1pt,
    every node/.style={
        circle,
        fill=black,
        inner sep=2.4pt
    },
    arr/.style={-{Triangle[length=2mm,width=1.5mm]}},
]
    % \node at (-7,-10) {\(F_3\)};
            
    \node[circle, fill=black, label=below:{\(u\)}]
            (r0) at (-4,-13) {};
    \node[circle, fill=black, label=below:{}]
            (r1) at (2,-13) {};
    \node[circle, fill=black, label=below:{}]
            (r2) at (8,-13) {};
    \node[circle, fill=black, label=below:{\(v\)}]
            (r3) at (14,-13) {};
    \node[circle, fill=black, label=below:{}]
            (r4) at (-3,-11) {};
    \node[circle, fill=black, label=below:{}]
            (r5) at (1,-11) {};
    \node[circle, fill=black, label=below:{}]
            (r6) at (3,-11) {};
    \node[circle, fill=black, label=below:{}]
            (r7) at (7,-11) {};
    \node[circle, fill=black, label=below:{}]
            (r8) at (9,-11) {};
    \node[circle, fill=black, label=below:{}]
            (r9) at (13,-11) {};
    \node[circle, fill=black, label=below:{}]
            (r10) at (3,-11) {};
    \node[circle, fill=black, label=below:{}]
            (r11) at (-2,-9) {};
    \node[circle, fill=black, label=below:{}]
            (r12) at (0,-9) {};
    \node[circle, fill=black, label=below:{}]
            (r13) at (4,-9) {};
    \node[circle, fill=black, label=below:{}]
            (r14) at (6,-9) {};
    \node[circle, fill=black, label=below:{}]
            (r15) at (10,-9) {};
    \node[circle, fill=black, label=below:{}]
            (r16) at (12,-9) {};
    \node[circle, fill=black, label=below:{}]
            (r17) at (-1,-7) {};
    \node[circle, fill=black, label=below:{}]
            (r18) at (5,-7) {};
    \node[circle, fill=black, label=below:{}]
            (r19) at (11,-7) {};

    \draw[arr] (r0) -- (r4);
    \draw[arr] (r1) -- (r5);
    \draw[arr] (r1) -- (r6);
    \draw[arr] (r2) -- (r7);
    \draw[arr] (r2) -- (r8);
    \draw[arr] (r3) -- (r9);
    \draw[arr] (r4) -- (r11);
    \draw[arr] (r5) -- (r12);
    \draw[arr] (r6) -- (r13);
    \draw[arr] (r7) -- (r14);
    \draw[arr] (r8) -- (r15);
    \draw[arr] (r9) -- (r16);
    \draw[arr] (r11) -- (r17);
    \draw[arr] (r12) -- (r17);
    \draw[arr] (r13) -- (r18);
    \draw[arr] (r14) -- (r18);
    \draw[arr] (r15) -- (r19);
    \draw[arr] (r16) -- (r19);

\end{tikzpicture}
  \end{column}
\end{columns}

\end{frame}

%---------------------------------------------------------


\begin{frame}{Stabilizacja płotków}
    \begin{figure}
        \centering
        \only<1-3>{\begin{tikzpicture}[
    scale=0.5,
    line width=1pt,
    every node/.style={
        circle,
        fill=black,
        inner sep=2.4pt
    },
    arr/.style={-{Triangle[length=2mm,width=1.5mm]}},
]

            
    \node[circle, fill=black, label=below:{\(u\)}]
            (q0) at (4,-4) {};
    \node[circle, fill=black, label=below:{}]
            (q1) at (8,-4) {};
    \node[circle, fill=black, label=below:{\(v\)}]
            (q2) at (12,-4) {};
    \node[circle, fill=black, label=below:{}]
            (q3) at (5,-2) {};
    \node[circle, fill=black, label=below:{}]
            (q4) at (7,-2) {};
    \node[circle, fill=black, label=below:{}]
            (q5) at (9,-2) {};
    \node[circle, fill=black, label=below:{}]
            (q6) at (11,-2) {};
    \node[circle, fill=black, label=left:{}]
        (q7) at (6,0) {};
    \node[circle, fill=black, label=left:{}]
        (q8) at (10,0) {}; 
    
        
    \draw[arr] (q0) -- (q3);
    \draw[arr] (q3) -- (q7);
    \draw[arr] (q4) -- (q7);
    \draw[arr] (q1) -- (q4);
    \draw[arr] (q1) -- (q5);
    \draw[arr] (q2) -- (q6);
    \draw[arr] (q5) -- (q8);
    \draw[arr] (q6) -- (q8);


\onslide<2->{
    \node[circle, fill=cfnavy, , label=below:{}]
        (q9) at (6,2) {};  
    \node[circle, fill=cfnavy, , label=below:{}]
        (q10) at (10,2) {}; 
    \draw[arr] (q7) -- (q9);
    \draw[arr] (q8) -- (q10);
}

\onslide<3>{
    \node[circle, fill=cfnavy, , label=below:{}]
        (q11) at (12,-2) {};  
    \node[circle, fill=cfnavy, , label=below:{}]
        (q12) at (12,0) {};
    \node[circle, fill=cfnavy, , label=below:{}]
        (q13) at (12,2) {};
    \draw[arr] (q2) -- (q11);
    \draw[arr] (q11) -- (q12);
    \draw[arr] (q12) -- (q13);
}
    
\end{tikzpicture}
}
        \only<4->{\begin{tikzpicture}[
    scale=0.5,
    line width=1pt,
    every node/.style={
        circle,
        fill=black,
        inner sep=2.4pt
    },
    arr/.style={-{Triangle[length=2mm,width=1.5mm]}},
]

            
    \node[circle, fill=black, label=below:{\(u\)}]
            (r0) at (-4,-13) {};
    \node[circle, fill=black, label=below:{}]
            (r1) at (2,-13) {};
    \node[circle, fill=black, label=below:{\(v\)}]
            (r2) at (8,-13) {};
    \node[circle, fill=black, label=below:{\(v\)}]
            (r3) at (14,-13) {};
    \node[circle, fill=black, label=below:{}]
            (r4) at (-3,-11) {};
    \node[circle, fill=black, label=below:{}]
            (r5) at (1,-11) {};
    \node[circle, fill=black, label=below:{}]
            (r6) at (3,-11) {};
    \node[circle, fill=black, label=below:{}]
            (r7) at (7,-11) {};
    \node[circle, fill=cfpink, label=right:{\(r_4\)}]
            (r8) at (9,-11) {};
    \node[circle, fill=cfpink, label=right:{\(r_4\)}]
            (r9) at (13,-11) {};
    \node[circle, fill=black, label=below:{}]
            (r10) at (3,-11) {};
    \node[circle, fill=cfgr, label=left:{\(r_1\)}]
            (r11) at (-2,-9) {};
    \node[circle, fill=cfgr, label=right:{\(r_1\)}]
            (r12) at (0,-9) {};
    \node[circle, fill=cfred, label=left:{\(r_2\)}]
            (r13) at (4,-9) {};
    \node[circle, fill=cfred, label=right:{\(r_2\)}]
            (r14) at (6,-9) {};
    \node[circle, fill=cfyel, label=left:{\(r_3\)}]
            (r15) at (10,-9) {};
    \node[circle, fill=cfyel, label=right:{\(r_3\)}]
            (r16) at (12,-9) {};
    \node[circle, fill=cfnavy, label=below:{}]
            (r17) at (-1,-7) {};
    \node[circle, fill=cfnavy, label=below:{}]
            (r18) at (5,-7) {};
    \node[circle, fill=cfnavy, label=below:{}]
            (r19) at (11,-7) {};

    \draw[arr] (r0) -- (r4);
    \draw[arr] (r1) -- (r5);
    \draw[arr] (r1) -- (r6);
    \draw[arr] (r2) -- (r7);
    \draw[arr] (r2) -- (r8);
    \draw[arr] (r3) -- (r9);
    \draw[arr] (r4) -- (r11);
    \draw[arr] (r5) -- (r12);
    \draw[arr] (r6) -- (r13);
    \draw[arr] (r7) -- (r14);
    \draw[arr] (r8) -- (r15);
    \draw[arr] (r9) -- (r16);
    \draw[arr] (r11) -- (r17);
    \draw[arr] (r12) -- (r17);
    \draw[arr] (r13) -- (r18);
    \draw[arr] (r14) -- (r18);
    \draw[arr] (r15) -- (r19);
    \draw[arr] (r16) -- (r19);

    % \draw[dashed] (r2) -- (r3) node[midway, above]{};
    % \draw[dashed] (r8) -- (r9) node[midway, above]{};
    % \draw[dashed] (r15) -- (r16) node[midway, above]{};
    % \draw[dashed] (r11) -- (r12) node[midway, above]{};
    % \draw[dashed] (r13) -- (r14) node[midway, above]{};
    
\end{tikzpicture}
}
    \end{figure}
    Jeśli istnieje płotek \(F_n(u,v)\), to istnieje także płotek \(F_{n+1}(u,v)\).
    \\
    \vspace{1em}
    \onslide<5>{
    Od pewnego momentu relacja się ustabilizuje: \vspace{-1em}
    \[ F_1 \subseteq F_2 \subseteq F_3 \subseteq \dots F_{m-1} \subsetneq F_m = F \]
    }
\end{frame}

%---------------------------------------------------------

\begin{frame}{al = 1}
    Graf ma długość algebraiczną 1, jeśli ustabilizowany płotek \(F\) łączy wszystkie wierzchołki.
    \[\mathrm{al}=1 \Longleftrightarrow \forall_{u,v\in G} \ F(u,v)\]
\end{frame}

%---------------------------------------------------------

\begin{frame}{Połączoność}
    Graf jest \emph{połączony} względem relacji \(R\), jeśli dowolną parę \(u, v \in V\) łączy ścieżka ,,góra dół''.
    \begin{figure}
        \centering
        \begin{tikzpicture}[
    scale=0.5,
    line width=1pt,
    every node/.style={
        circle,
        fill=black,
        inner sep=2.4pt
    },
    arr/.style={-{Triangle[length=2mm,width=1.5mm]}},
]
   % \node at (-4,-2) {\(F_2\)};
            
    \node[circle, fill=black, label=below:{\(u\)}]
            (q0) at (4,-4) {};
    \node[circle, fill=black, label=below:{}]
            (q1) at (8,-4) {};
    \node[circle, fill=black, label=below:{}]
            (q2) at (12,-4) {};
    \node[circle, fill=black, label=below:{}]
        (q7) at (6,0) {};
    \node[circle, fill=black, label=below:{}]
        (q8) at (10,0) {};   

    \node[circle, fill=black, inner sep=1pt] at (14,-2) {};
    \node[circle, fill=black, inner sep=1pt] at (15,-2) {};
    \node[circle, fill=black, inner sep=1pt] at (16,-2) {};

    \node[circle, fill=black, label=below:{}]
            (qa) at (18,-4) {};
    \node[circle, fill=black, label=below:{\(v\)}]
            (qb) at (22,-4) {};
    \node[circle, fill=black, label=below:{}]
        (qc) at (20,0) {};   
    \draw[arr, cfred] (q0)  -- node[midway,left,
                 draw=none,
                 fill=none,
                 circle=false]  {\(R\)} (q7);
    \draw[arr, cfred] (q1)  --  (q7);
    \draw[arr, cfred] (q1)  --  (q8);
    \draw[arr, cfred] (q2)  -- (q8);
    \draw[arr, cfred] (qa)  -- node[midway,left,
                 draw=none,
                 fill=none,
                 circle=false] {\(R\)} (qc);
    \draw[arr, cfred] (qb)  --  (qc);
\end{tikzpicture}
    \end{figure}
\end{frame}
